\documentclass[../Main.tex]{subfiles}
\begin{document}
\chapter{2017-2018学年微积分（一）（上）期中考试}

\section{基本计算题(每小题 6 分，共 60 分)}
\begin{enumerate}
    \item 设数列$x_{n}$满足$x_{1}=1,x_{n+1}=\sin x_{n}(n>1)$，求极限$\lim_{n\to\infty}
              x_{n}$.
          \vspace{11em}

    \item 求数列极限$\lim_{n\to \infty}\cos\left(\pi\sqrt{n^{2}+n}\right)$.
          \vspace{10em}

    \item 计算极限$l=\lim_{x\to 0}\frac{x\cos x-\sin x}{(\cos x-1)\ln(1+x)}$.
          \vspace{11em}

    \item 已知$\lim_{x\to1}\frac{ax^{3}+2x+1}{x-1}=b$，求常数$a,b$的值.
          \vspace{10em}

    \item 求极限$l=\lim_{x\to1}\left[\frac{2}{1-\mathrm{e}^{\frac{x}{1-x}}}-\frac{|\sin(x-1)|}{x-1}
                  \right]$.
          \vspace{10em}

    \item 指出函数$f(x)=\frac{|x|}{\tan x}$的间断点，并确定间断点的类型.
          \vspace{10em}

    \item 设函数$y=(2+x)^{\sin x}+\frac{1}{x+1}(x>-2)$，求微分$\left.\mathrm{d}y\right
              |_{x=0}$.
          \vspace{11em}

    \item 设函数$v=f(u)$有反函数$u=\varphi(v)$，满足$f(0)=\frac{\pi}{2}$，且$\varphi
              (v)$是可导的，在$v=0$的某个邻域内有$\varphi'(v)=\frac{1}{2+\sin v}$，求复合函数$y
              =f^{2}(x)$在$x=0$处的导数.
          \vspace{12em}

    \item 设$y=(x-1)\ln x$，求$y^{(10)}(1)$.
          \vspace{8em}

    \item 设函数$y=y(x)$由方程$\begin{cases}
                  x=t\cos t, \\
                  y=t\sin t
              \end{cases}$确定，求在$t=0$时的导数$\frac{\mathrm{d}y}{\mathrm{d}x}$和$\frac{\mathrm{d}^{2}y}{\mathrm{d}x^{2}}$.
          \vspace{8em}
\end{enumerate}

\section{综合题(每小题 6 分，共 30 分)}
\begin{enumerate}
    \setcounter{enumi}{10}
    \item 设函数$f(x)$在原点附近有界，$F(x)=f(x)\cdot\sin\left(x^{2}\right)$，计算导数$F
              '(0)$.
          \vspace{8em}

    \item 求函数$f(x)=
              \begin{cases}
                  \frac{1-\cos x}{x}, & x\neq0, \\
                  0,                  & x=0
              \end{cases}$的导函数$f'(x)$，并讨论$f'(x)$的连续性.
          \vspace{8em}

    \item 设函数$f(x)$在$x=0$的某个邻域内有界，满足$\lim_{x\to0}\frac{\sqrt[n]{1+f(x)\sin
                      x}-1}{\mathrm{e}^{3x}-1}=2$，其中$n$为正整数.求$\lim_{x\to0}f(x)$.
          \vspace{9em}

    \item 求无穷小量$u(x)=x-\arctan x(x\to0)$的主部和阶数.
          \vspace{10em}

    \item 一个长方体的铁皮盒子，其对角线的长度随着长宽高的变化而连续变化.当长宽高分别是$3
              \mathrm{m},4\mathrm{m},5\mathrm{m}$时，如果此时对角线长度增加的速率为$5\sqrt{2}
              \mathrm{m/s}$，长宽增加的速率分别为$8\mathrm{m/s}$和$9\mathrm{m/s}$，问此时高是在增加还是在减少？增加或减少的速率为多少？
          \vspace{12em}
\end{enumerate}

\section{证明题(每小题 5 分，共 10 分)}
\begin{enumerate}
    \setcounter{enumi}{15}
    \item 证明函数$f(x)=
              \begin{cases}
                  x^{2}, & x\in\mathbf{Q},                   \\
                  0,     & x\in\mathbf{R}\setminus\mathbf{Q}
              \end{cases}$在$x=0$处可导，但函数本身除零点外处处不连续.
          \vspace{12em}

    \item 设$f(x)$在$[0,3]$上连续，在$(0,3)$内一阶可导，且$f(0)=0,f(1)=3,f(3
              )=1$，证明至少存在一点$\xi\in(0,3)$使得$f'(\xi)=0$.
\end{enumerate}

\chapter{2017-2018学年微积分（一）（上）期中考试参考答案}
\section{基本计算题(每小题 6 分，共 60 分)}
\begin{enumerate}
    \item \textbf{Solution}. 显然$0<x_{2}=\sin x_{1}<\frac{\pi}{2}$.设$0<x_{k}<\frac{\pi}{2}$，则$0
              <x_{k+1}=\sin x_{k}<x_{k}<\frac{\pi}{2}$，

          由数学归纳法知$\{x_{n}\}$单调递减且有下界0，故$\{x_{n}\}$收敛.

          设$\lim_{n\to\infty}x_{n}=l$，在方程$x_{n+1}=\sin x_{n}$两边取极限得到$l=\sin
              l$，解得$l=0$.

    \item \textbf{Solution}.
          \[
              \begin{aligned}
                  \lim_{n\to\infty}\cos\left(\pi\sqrt{n^2+n}\right)= & \lim_{n\to\infty}(-1)^{n}\cos\left(\pi\sqrt{n^2+n}-n\pi\right)               \\
                  =                                                  & \lim_{n\to\infty}(-1)^{n}\cos\left(\frac{n\pi}{\sqrt{n^2+n}+n}\right)        \\
                  =                                                  & \lim_{n\to\infty}(-1)^{n}\cos\left(\frac{\pi}{\sqrt{1+\frac{1}{n}}+1}\right) \\
                  =                                                  & 0.
              \end{aligned}
          \]

    \item \textbf{Solution}.
          \[
              \begin{aligned}
                  l= & \lim_{x\to0}\frac{x\cos x-\sin x}{(\cos x-1)\ln(1+x)}= \lim_{x\to0}\frac{x\left(1-\frac{x^2}{2}+o(x^2)\right)-\left(x-\frac{x^3}{6}+o(x^3)\right)}{\left(1-\frac{x^2}{2}+o(x^2)-1\right)\left(x+o(x)\right)} \\
                  =  & \lim_{x\to0}\frac{-\frac{1}{3}x^3+o(x^3)}{-\frac{1}{2}x^3+o(x^3)}                                                                                                                                            \\
                  =  & \frac{2}{3}.
              \end{aligned}
          \]

    \item \textbf{Solution}.由
          \[
              \lim_{x\to1}(ax^{3}+2x+1)=a+3=\lim_{x\to1}\frac{ax^{3}+2x+1}{x-1}\cdot(x-
              1)=0
          \]

          解得$a=-3$，从而
          \[
              b=\lim_{x\to1}\frac{-3x^{3}+2x+1}{x-1}=\lim_{x\to1}\frac{-9x^{2}+2}{1}= -
              7.
          \]

    \item \textbf{Solution}. 当$x\to1^{-}$时，$\mathrm{e}^{\frac{x}{1-x}}\to+\infty$，$\frac{|\sin(x-1)|}{x-1}
              \to-1$，所以
          \[
              \lim_{x\to1^-}\left[\frac{2}{1-\mathrm{e}^{\frac{x}{1-x}}}-\frac{|\sin(x-1)|}{x-1}
                  \right]= 0-(-1)=1.
          \]

          当$x\to1^{+}$时，$\mathrm{e}^{\frac{x}{1-x}}\to0$，$\frac{|\sin(x-1)|}{x-1}
              \to1$，所以
          \[
              \lim_{x\to1^+}\left[\frac{2}{1-\mathrm{e}^{\frac{x}{1-x}}}-\frac{|\sin(x-1)|}{x-1}
                  \right]= 2-1=1.
          \]

          因此$l=1$.

    \item \textbf{Solution}. 函数$f(x)$的间断点为$x=k\pi,k\pi+\frac{\pi}{2}(k\in\mathbf{Z}
              )$.

          当$x=0$时，$\lim_{x\to0^+}\frac{|x|}{\tan x}=1$，$\lim_{x\to0^-}\frac{|x|}{\tan
                  x}=-1$，所以$x=0$为跳跃间断点.

          当$x=k\pi(k\neq0)$时，$\lim_{x\to k\pi}\frac{|x|}{\tan x}=\infty$，所以$x=k
              \pi(k\neq0)$为无穷间断点.

          当$x=k\pi+\frac{\pi}{2}$时，$\lim_{x\to k\pi+\frac{\pi}{2}}\frac{|x|}{\tan
                  x}=0$，所以$x=k\pi+\frac{\pi}{2}$为可去间断点.

    \item \textbf{Solution}. 记$u=(2+x)^{\sin x}$，$v=\frac{1}{x+1}$，则$\ln u = \sin
              x\cdot\ln(2+x)$，所以
          \[
              \begin{gathered}
                  \mathrm{d}u =u\cdot\,\mathrm{d}\left(\sin x\ln(2+x)\right)=(2+x)^{\sin x}\left(\cos
                  x\ln(2+x)+\frac{\sin x}{2+x}\right)\,\mathrm{d}x,\\
                  \mathrm{d}v =-\frac{1}{(x+1)^{2}}\,\mathrm{d}x.
              \end{gathered}
          \]

          代入$x=0$得$\left.\mathrm{d}y\right|_{x=0}=\left.\mathrm{d}u\right|_{x=0}+\left
              .\mathrm{d}v\right|_{x=0}=(\ln2-1)\,\mathrm{d}x$.

    \item \textbf{Solution}.
          \[
              \begin{aligned}
                  y'(0) & =\left.2f(x)f'(x)\right|_{x=0}                          \\
                        & = 2f(0)f'(0)                                            \\
                        & =\frac{2f(0)}{\varphi'\left(\frac{\pi}{2}\right)}=3\pi.
              \end{aligned}
          \]

    \item \textbf{Solution}. $y'=\ln x+1-\frac{1}{x}$，所以
          \[
              \begin{aligned}
                  y^{(10)}= & \left(\ln x\right)^{(9)}-\left(\frac{1}{x}\right)^{(9)} \\
                  =         & \frac{(-1)^8\,8!}{x^9}-\frac{(-1)^9\,9!}{x^{10}}.
              \end{aligned}
          \]

          代入$x=1$得$y^{(10)}(1)=8!+9!=10\cdot8!$.

    \item \textbf{Solution}. $\frac{\mathrm{d}y}{\mathrm{d}x}=\frac{\frac{\mathrm{d}y}{\mathrm{d}t}}{\frac{\mathrm{d}x}{\mathrm{d}t}}
              =\frac{\sin t+t\cos t}{\cos t-t\sin t}$，所以$\left.\frac{\mathrm{d}y}{\mathrm{d}x}
              \right|_{t=0}=0$.

          $\frac{\mathrm{d}^{2}y}{\mathrm{d}x^{2}}=\frac{\mathrm{d}}{\mathrm{d}t}\left
              (\frac{\mathrm{d}y}{\mathrm{d}x}\right)\cdot\frac{\mathrm{d}t}{\mathrm{d}x}
              =\left(\frac{\sin t+t\cos t}{\cos t-t\sin t}\right)'_{t}\cdot\frac{1}{\cos
                  t-t\sin t}=\frac{2+t^{2}}{(\cos t-t\sin t)^{3}}$，所以$\left.\frac{\mathrm{d}^{2}y}{\mathrm{d}x^{2}}
              \right|_{t=0}=2$.
\end{enumerate}

\section{综合题(每小题 6 分，共 30 分)}
\begin{enumerate}
    \setcounter{enumi}{10}
    \item \textbf{Solution}.
          \[
              \begin{aligned}
                  F'(0)= & \lim_{x\to0}\frac{f(x)\sin(x^2)-0}{x}              \\
                  =      & \lim_{x\to0}\frac{x^2f(x)}{x}=\lim_{x\to0}xf(x)=0.
              \end{aligned}
          \]

    \item \textbf{Solution}. 当$x\neq0$时，$f'(x)=\frac{x\sin x-1+\cos x}{x^{2}}$；

          当$x=0$时，$f'(0)=\lim_{x\to0}\frac{\frac{1-\cos x}{x}-0}{x}=\lim_{x\to0}\frac{1-\cos
                  x}{x^{2}}=\frac{1}{2}$.

          所以$f'(x)=
              \begin{cases}
                  \frac{x\sin x-1+\cos x}{x^2}, & x\neq0, \\
                  \frac{1}{2},                  & x=0.
              \end{cases}$

          当$x\neq0$时，$f'(x)$显然连续.又
          \[
              \lim_{x\to0}f'(x)=\lim_{x\to0}\frac{x\sin x-1+\cos x}{x^{2}}=1-\frac{1}{2}
              =\frac{1}{2}=f'(0),
          \]

          所以$f'(x)$在$x=0$处连续.综上所述，$f'(x)$在$\mathbf{R}$上处处连续.

    \item \textbf{Solution}. 由于$f(x)$在$x=0$的某个邻域内有界，所以$\lim_{x\to0}
              f(x)\sin x=0$，故
          \[
              \begin{aligned}
                  \lim_{x\to0}\frac{\sqrt[n]{1+f(x)\sin x}-1}{\mathrm{e}^{3x}-1} & =\lim_{x\to0}\frac{f(x)\sin x}{n\cdot 3x} \\
                                                                                 & = \frac{1}{3n}\lim_{x\to0}f(x)=2,
              \end{aligned}
          \]

          从而$\lim_{x\to0}f(x)=6n$.

    \item \textbf{Solution}. 设$\lim_{x\to0}\frac{u(x)}{cx^{r}}=1$，则
          \[
              1=\lim_{x\to0}\frac{x-\arctan x}{cx^{r}}=\lim_{x\to0}\frac{1-\frac{1}{1+x^{2}}}{crx^{r-1}}
              =\lim_{x\to0}\frac{1}{1+x^{2}}\cdot\frac{x^{2}}{crx^{r-1}},
          \]

          因此必须$\begin{cases}
                  r-1=2, \\
                  cr=1
              \end{cases}$，解得$r=3,c=\frac{1}{3}$.

          所以$u(x)$的主部为$\frac{1}{3}x^{3}$，阶数为3.

    \item \textbf{Solution}. 设$t$时刻长方体的长宽高及对角线长分别为$x(t),y(
              t),z(t),s(t)$，则
          \[
              s^{2}(t)=x^{2}(t)+y^{2}(t)+z^{2}(t).
          \]

          方程两边对$t$求导得
          \[
              s(t)s'(t)=x(t)x'(t)+y(t)y'(t)+z(t)z'(t).
          \]

          将$x(t)=3\mathrm{m},y(t)=4\mathrm{m},z(t)=5\mathrm{m},s(t)=5\sqrt{2}\mathrm{m}
              ,x'(t)=8\mathrm{m/s},y'(t)=9\mathrm{m/s},s'(t)=5\sqrt{2}\mathrm{m/s}$

          代入上式得$z'(t)=-2\mathrm{m/s}$，说明此时长方体的高在减少，减少的速率为$2\mathrm{m/s}$.
\end{enumerate}

\section{证明题(每小题 5 分，共 10 分)}
\begin{enumerate}
    \setcounter{enumi}{15}
    \item \textbf{Proof}. 计算$\lim_{x\to0}\frac{f(x)-0}{x-0}=
              \begin{cases}
                  \lim_{x\to0}\frac{x^2}{x}=0, & x\in\mathbf{Q},                    \\
                  0,                           & x\in\mathbf{R}\setminus\mathbf{Q}.
              \end{cases}$

          故$f(x)$在$x=0$处可导，且$f'(0)=0$.

          考虑非零点$a$. 若$a\in\mathbf{Q}$，取点列$\{x_{n},x_{n}\in\mathbf{R}\setminus
              \mathbf{Q}\}$使得$\lim_{n\to\infty}x_{n}=a$，

          则$\lim_{n\to\infty}f(x_{n})=0\neq f(a)=a^{2}$，说明$f(x)$在$x=a$处不连续.

          同理可证当$a\in\mathbf{R}\setminus\mathbf{Q}$时，$f(x)$也不连续.所以函数除零点外处处不连续.

    \item \textbf{Proof}. 由于$f(x)$在$[0,3]$上连续，由介值定理，存在$\eta\in
              (0,1)$使得$f(0)<f(\eta)=1<f(1)$.

          $f(x)$在$[\eta,3]$上连续，在$(\eta,3)$内可导，由Rolle定理知存在$\xi\in(\eta
              ,3)\subset(0,3)$使得$f'(\xi)=0$.
\end{enumerate}
\end{document}
